\chapter{CƠ SỞ LÝ THUYẾT VÀ CÁC CÔNG TRÌNH LIÊN QUAN}

\section{Cơ chế gọi công cụ (Tool Calling / Function Calling) trong Mô hình Ngôn ngữ}

Các công trình về mô hình ngôn ngữ sử dụng công cụ phát triển theo nhiều hướng. Toolformer [2] khảo sát cách mô hình tự học vị trí gọi API; Gorilla [3] tập trung vào khả năng sử dụng số lượng lớn API; Berkeley Function-Calling Leaderboard (BFCL) [4] cung cấp các tình huống và thước đo đánh giá lời gọi hàm. Các công trình này đặt nền tảng cho việc phân tích cả lựa chọn công cụ lẫn tham số trong khóa luận.

Gọi $X = (x_1, x_2, \dots, x_{|X|})$ là truy vấn và $\mathcal{T} = \{t_1, t_2, \dots, t_N\}$ là danh mục $N$ công cụ. Mỗi công cụ $t_i = (n_i, d_i, \mathcal{S}_i)$ có tên $n_i$, mô tả $d_i$ và lược đồ tham số $\mathcal{S}_i$ theo JSON Schema. Hệ thống ánh xạ $(X, \mathcal{T})$ thành chuỗi lời gọi $\mathcal{Y} = [(t^*_1, \mathcal{A}_1), \dots, (t^*_K, \mathcal{A}_K)]$, với $K=0$ khi không cần gọi công cụ. Ở mỗi lời gọi, $t^*_k \in \mathcal{T}$ và $\mathcal{A}_k = \{(p_{k,j}, v_{k,j})\}_{j=1}^{M_k}$ là các cặp tên–giá trị tham số phù hợp với $\mathcal{S}_{t^*_k}$.

Trong cách tiếp cận tạo sinh, mô tả các công cụ khả dụng được đưa vào ngữ cảnh đầu vào để mô hình sinh lời gọi theo phân phối có điều kiện:
$$P(\mathcal{Y} \mid X, \mathcal{T}) = \prod_{t=1}^{|\mathcal{Y}|} P(y_t \mid y_{<t}, X, \mathcal{T})$$
Khi số công cụ tăng, việc đưa toàn bộ schema vào prompt làm dài chuỗi đầu vào, tăng nhu cầu bộ nhớ và có thể gây khó khăn cho việc phân biệt các công cụ tương đồng. Chi phí chú ý toàn phần tăng theo $\mathcal{O}(L^2)$ với độ dài chuỗi $L$, nhưng không đại diện cho toàn bộ các lớp của kiến trúc Qwen3.5 [17]. Tác động thực tế lên độ chính xác và bộ nhớ được kiểm tra bằng stress test ở Chương 5.

\section{Các công trình và bộ chuẩn đánh giá Tool Calling trong và ngoài nước}

Các tập dữ liệu huấn luyện và đánh giá công cụ có nguồn gốc khác nhau. Glaive Function Calling v2 [7] cung cấp các lượt đối thoại tổng hợp có lời gọi hàm và phản hồi trò chuyện thông thường; khóa luận sử dụng nguồn này sau các bước lọc, chuẩn hóa và dịch thuật được trình bày ở Chương 3.

Salesforce công bố bộ dữ liệu xLAM-60k được tạo bằng quy trình APIGen, trong đó có các truy vấn gọi nhiều hàm và các kiểu tham số phức tạp [11]. ToolLLM và ToolBench [5] khai thác danh mục API lớn cho huấn luyện và đánh giá; BFCL [4] đánh giá các năng lực gọi hàm qua nhiều dạng tình huống. Các nguồn này làm rõ bài toán lựa chọn công cụ và điền tham số, nhưng khác nhau về tập công cụ, cấu trúc truy vấn và cách tính độ đo; vì vậy, điểm số giữa các bộ chuẩn không thể đối chiếu trực tiếp.

Phần lớn các nguồn dữ liệu trên thiên về tiếng Anh. Ersoy và cộng sự [1] khảo sát dịch dữ liệu gọi công cụ sang tiếng Ả Rập và tinh chỉnh mô hình nhỏ, cho thấy hướng bản địa hóa cần được kiểm chứng trên từng ngôn ngữ thay vì suy từ kết quả tiếng Anh.

Đối với tiếng Việt, bộ \textit{Vietnamese Function Calling Benchmark} do phamhai công bố [23] gồm 2,899 truy vấn đơn lượt trên 159 hàm thuộc nhiều lĩnh vực và báo cáo cả độ chính xác tên hàm lẫn độ khớp toàn bộ lời gọi. Nguồn này cung cấp một điểm tham chiếu cho đánh giá gọi công cụ tiếng Việt. Tuy nhiên, tập kiểm thử và giao thức của nguồn này khác với các bộ chuẩn trong khóa luận; các điểm số đã công bố không được dùng để so sánh trực tiếp. Thẻ dữ liệu không trình bày phép đối chiếu giữa kiến trúc tạo sinh và kiến trúc phân tách dưới cùng một giao thức, cũng không báo cáo stress test khi số công cụ tăng đến hàng nghìn.

Từ các công trình được khảo cứu, khóa luận tập trung vào ba khoảng trống trong phạm vi thực nghiệm của mình: đánh giá cùng một lược đồ dữ liệu cho hai kiến trúc gọi công cụ tiếng Việt; kiểm tra khả năng từ chối và xử lý công cụ chưa từng gặp trên dữ liệu bản địa; và đo độ bền khi danh mục công cụ mở rộng từ $N=3$ đến $N=1{,}000$. Các phép đối chiếu được giới hạn trong bộ dữ liệu và giao thức đã thiết lập, không hàm ý rằng nghiên cứu đã bao quát mọi công trình tiếng Việt hiện có.

\section{Mô hình Ngôn ngữ Nhỏ (SLMs) và Kỹ thuật Tinh chỉnh Tham số Hiệu quả (PEFT/QLoRA)}

Trong bối cảnh bài toán triển khai thực tế đòi hỏi tối ưu hóa tài nguyên phần cứng và bảo vệ dữ liệu doanh nghiệp, các Mô hình Ngôn ngữ Nhỏ (Small Language Models - SLMs) với số lượng tham số từ 1 tỷ đến 4 tỷ là một hướng tiếp cận đáng khảo sát. Khóa luận sử dụng Qwen3.5 ở hai cấu hình 2B và 4B; năng lực của các cấu hình này được đánh giá trực tiếp trên benchmark thay vì suy ra chỉ từ quy mô tiền huấn luyện.

Theo cấu hình công bố của Qwen3.5-2B và Qwen3.5-4B [17], mô hình xen kẽ các lớp linear attention và full attention. Ở các lớp full attention, số đầu truy vấn và số đầu khóa–giá trị khác nhau, phù hợp với cơ chế Grouped-Query Attention (GQA) [13]; cấu hình cũng khai báo tham số Rotary Position Embedding (RoPE) [14]. Vì vậy, không nên áp dụng độ phức tạp của full attention cho toàn bộ các lớp. RoPE biểu diễn thông tin vị trí qua phép quay trên các chiều đặc trưng:
$$\mathbf{R}_{\Theta, m}^d = \text{diag}\left(\mathbf{R}_{\theta_1, m}, \mathbf{R}_{\theta_2, m}, \dots, \mathbf{R}_{\theta_{d/2}, m}\right)$$
Đối với mạng truyền thẳng, SwiGLU là một biến thể cổng dùng hàm SiLU/Swish [15]. Cấu hình text của Qwen3.5 khai báo hàm kích hoạt SiLU [17]; biểu thức dưới đây mô tả họ cơ chế cổng, không tự nó xác nhận tốc độ hội tụ của các checkpoint trong nghiên cứu:
$$\text{SwiGLU}(x) = \left(x W_{gate} \cdot \sigma(x W_{gate})\right) \otimes (x W_{up})$$

Để tinh chỉnh các mô hình SLM trong giới hạn bộ nhớ, nghiên cứu áp dụng QLoRA [6] qua Unsloth [16]. QLoRA kết hợp lượng tử hóa NF4, lượng tử hóa kép và tối ưu hóa bộ nhớ; trọng số gốc $W_0 \in \mathbb{R}^{d_{out} \times d_{in}}$ được giữ cố định ở dạng lượng tử hóa, còn các ma trận thích ứng hạng thấp $A \in \mathbb{R}^{r \times d_{in}}$ và $B \in \mathbb{R}^{d_{out} \times r}$ được cập nhật:
$$W = W_0 + \Delta W = W_0 + \frac{\alpha}{r} (B \cdot A)$$
với hạng ma trận $r \ll \min(d_{in}, d_{out})$ và hệ số co giãn siêu tham số $\alpha$. Cơ chế này làm giảm số lượng tham số cần cập nhật so với tinh chỉnh toàn bộ; mức tiết kiệm bộ nhớ và khả năng chạy trên một GPU cụ thể phụ thuộc vào kích thước mô hình, độ dài ngữ cảnh và cấu hình huấn luyện.

\section{Kiến trúc Phân tách: Truy hồi Ngữ nghĩa Dày (Bi-Encoder) và Trích xuất Tham số (Cross-Encoder)}

Nhằm giảm ảnh hưởng của quá tải ngữ cảnh và độ trễ sinh từ tự hồi quy, hướng tiếp cận phân tách (Method 2) chia bài toán Tool Calling thành hai khối thành phần chuyên biệt và xử lý tuần tự theo mô hình suy luận phân biệt (Discriminative Pipeline):

\begin{verbatim}
graph TD
    classDef io fill:#e1f5fe,stroke:#0288d1,stroke-width:2px,color:#01579b;
    classDef model fill:#e8f5e9,stroke:#388e3c,stroke-width:2px,color:#1b5e20;
    classDef inter fill:#fff8e1,stroke:#ffa000,stroke-width:2px,color:#e65100;
    classDef head fill:#ede7f6,stroke:#512da8,stroke-width:2px,color:#311b92;

    Query["Câu truy vấn của người dùng (User Query)"]:::io
    ToolPool["Kho danh mục công cụ (Tool Pool: N APIs)"]:::io

    subgraph Stage1 ["Giai đoạn 1: Lựa chọn công cụ (Tool Retrieval)"]
        BiEncoder["Mô hình Bi-Encoder (BGE-M3 + LoRA)"]:::model
        CosineSim["Xếp hạng tương đồng Cosine & Lọc ngưỡng tau, delta"]:::inter
    end

    subgraph Stage2 ["Giai đoạn 2: Trích xuất tham số theo lược đồ (Cross-Encoder Extraction)"]
        CrossEncoder["Mô hình Cross-Encoder (XLM-RoBERTa Backbone)"]:::model
        subgraph HierarchicalHeads ["Các đầu phân loại phân cấp (Hierarchical Heads)"]
            HasValue["has_value Head (Phân loại nhị phân Có/Không)"]:::head
            SpanHead["span_head (Dự đoán vị trí Start / End)"]:::head
            EnumHead["enum_head (Phân loại danh mục tĩnh)"]:::head
            BoolHead["bool_head (Phân loại giá trị True / False)"]:::head
        end
        Normalizer["Module Chuẩn hóa Giá trị (Vietnamese Value Normalizer)"]:::inter
    end

    Output["Lệnh gọi hàm hoàn chỉnh (JSON Function Call)"]:::io

    Query --> BiEncoder
    ToolPool --> BiEncoder
    BiEncoder --> CosineSim
    CosineSim -->|"Top-K công cụ ứng viên"| CrossEncoder
    Query -->|"Ngữ cảnh truy vấn (Context)"| CrossEncoder

    CrossEncoder --> HasValue
    CrossEncoder --> SpanHead
    CrossEncoder --> EnumHead
    CrossEncoder --> BoolHead

    HasValue -->|"Tham số có mặt"| Normalizer
    SpanHead --> Normalizer
    EnumHead --> Normalizer
    BoolHead --> Normalizer

    Normalizer --> Output
\end{verbatim}

\textit{Hình 2.1: Sơ đồ nguyên lý kiến trúc phân tách hai giai đoạn (Method 2).}


\section{1. Giai đoạn 1: Truy hồi công cụ ngữ nghĩa dày (Bi-Encoder Retrieval)}

Giai đoạn đầu định vị các công cụ $t^* \in \mathcal{T}$ có khả năng đáp ứng truy vấn. Nhóm nghiên cứu dùng Bi-Encoder dựa trên BGE-M3 [10] để mã hóa truy vấn $q$ và mô tả công cụ $t$ thành các vector nhúng $u, v \in \mathbb{R}^d$ thông qua phép gom trung bình có trọng số (Mean Pooling):
$$u = \text{BiEncoder}(q), \quad v = \text{BiEncoder}(t)$$
Mức độ phù hợp giữa truy vấn và công cụ được xác định qua hàm khoảng cách Cosine Similarity:
$$s(q, t) = \frac{u^\top v}{\|u\|_2 \|v\|_2}$$

Để học biểu diễn ngữ nghĩa của công cụ tiếng Việt, mô hình sử dụng Multiple Negatives Ranking Loss (MNRL) kết hợp kỹ thuật đệm gradient cho lô lớn [22]. Hàm mất mát cho một lô gồm $B$ cặp dương $(q_i, t_i^+)$ được định nghĩa như sau:
$$\mathcal{L}_{\text{MNRL}} = -\frac{1}{B} \sum_{i=1}^B \log \frac{\exp\left(s(q_i, t_i^+) / \tau\right)}{\exp\left(s(q_i, t_i^+) / \tau\right) + \sum_{j \neq i} \exp\left(s(q_i, t_j^+) / \tau\right) + \sum_{k=1}^{M} \exp\left(s(q_i, n_{i,k}^-) / \tau\right)}$$
trong đó $\tau$ là siêu tham số nhiệt độ (temperature scaling), $t_j^+$ là các mẫu âm ngẫu nhiên nội lô (in-batch negatives), và $n_{i,k}^-$ là các mẫu âm tính khó (hard negatives) được khai phá độc lập thông qua mô hình giáo viên ở vòng huấn luyện đầu tiên. Chiến lược này buộc mô hình phải phân biệt ranh giới ngữ nghĩa cực kỳ tinh tế giữa các công cụ có chức năng tương đồng trong cùng một lĩnh vực nghiệp vụ.

\section{2. Giai đoạn 2: Trích xuất tham số theo lược đồ (Cross-Encoder Extraction)}

Sau khi truy xuất Top-$K$ công cụ ứng viên, giai đoạn thứ hai trích xuất giá trị cho từng thuộc tính trong lược đồ $\mathcal{S}$. Thiết kế tham khảo bài toán hỏi đáp trích xuất trên BERT [8] và sử dụng backbone đa ngôn ngữ XLM-RoBERTa [9]. Với mỗi tham số $p_j$, đầu vào ghép truy vấn người dùng làm ngữ cảnh (Context) và mô tả tham số làm câu hỏi (Question):
$$\mathbf{X}_{\text{input}} = \text{[CLS]} \circ \text{Query} \circ \text{[SEP]} \circ \text{Parameter Prompt}(p_j, \mathcal{S}) \circ \text{[SEP]}$$

Biểu diễn ngữ cảnh tương tác sâu giữa truy vấn và câu hỏi tại đầu ra của mô hình được định tuyến qua hệ thống các đầu phân loại phân cấp chuyên biệt. Đầu phân loại nhị phân hiện diện (\texttt{has_value_head}) áp dụng hàm kích hoạt Sigmoid lên biểu diễn ẩn của token \texttt{[CLS]} để xác định xem tham số $p_j$ có thực sự xuất hiện và cần được gán giá trị trong truy vấn hay không:
$$\hat{y}_{\text{exist}} = \sigma\left(\mathbf{w}_{\text{exist}}^\top \mathbf{h}_{\text{[CLS]}} + b_{\text{exist}}\right)$$
với hàm mất mát tương ứng là entropy chéo nhị phân $\mathcal{L}_{\text{exist}} = \text{BCE}(\hat{y}_{\text{exist}}, y_{\text{exist}})$.

Khi tham số được xác nhận tồn tại, biểu diễn ẩn được định tuyến tới các nhánh chuyên biệt tùy theo kiểu dữ liệu. Đối với các tham số dạng chuỗi tự do, số thực hoặc số nguyên, đầu trích xuất đoạn văn bản (\texttt{span_head}) dự đoán phân phối xác suất của vị trí bắt đầu $i$ và kết thúc $j$ trên chuỗi token truy vấn:
$$P_{\text{start}}(i) = \frac{\exp(\mathbf{w}_s^\top \mathbf{h}_i)}{\sum_k \exp(\mathbf{w}_s^\top \mathbf{h}_k)}, \quad P_{\text{end}}(j) = \frac{\exp(\mathbf{w}_e^\top \mathbf{h}_j)}{\sum_k \exp(\mathbf{w}_e^\top \mathbf{h}_k)}$$
Hàm mất mát span là tổng entropy chéo phân loại: $\mathcal{L}_{\text{span}} = \text{CE}(P_{\text{start}}, y_s) + \text{CE}(P_{\text{end}}, y_e)$. Với các tham số có miền giá trị hữu hạn, \texttt{enum_head} và \texttt{bool_head} dự đoán nhãn trong danh mục schema. Cơ chế này giới hạn định dạng và miền giá trị có thể xuất ra, nhưng vẫn có thể chọn sai nhãn.

Tổng hàm mất mát liên hợp của mô hình Cross-Encoder là sự kết hợp có trọng số của các mục tiêu:
$$\mathcal{L}_{\text{total}} = \lambda_1 \mathcal{L}_{\text{exist}} + \mathbb{I}(y_{\text{exist}}=1) \cdot \left[ \lambda_2 \mathcal{L}_{\text{span}} + \lambda_3 \mathcal{L}_{\text{enum}} + \lambda_4 \mathcal{L}_{\text{bool}} \right]$$
Kiến trúc này cho phép suy luận độc lập cho từng tham số; bước hậu xử lý theo schema giúp đầu ra tuân thủ ràng buộc cú pháp trong protocol structured-output được sử dụng.

\section{Những thách thức đặc thù của ngôn ngữ tiếng Việt trong bài toán Tool Calling}

Bài toán gọi công cụ tiếng Việt cần xử lý đồng thời ranh giới thực thể, cách diễn đạt khẩu ngữ và biến thể chính tả.

Tiếng Việt là ngôn ngữ đơn lập; khoảng trắng thường ngăn cách âm tiết nhưng không luôn đánh dấu ranh giới từ. Chẳng hạn, “máy tính xách tay” gồm nhiều âm tiết tạo thành một đơn vị nghĩa. Khi kết hợp với tách từ phụ (subword tokenization), đặc điểm này có thể làm mô hình dự đoán span cắt cụt hoặc lấy thừa phần biên của giá trị tham số.

Các biểu thức như “hai củ rưỡi” ($2{,}500{,}000$ đồng), “năm xị” ($500{,}000$ đồng) hay “thứ Năm tuần tới” đòi hỏi chuyển đổi từ chuỗi bề mặt sang giá trị phù hợp với schema. Nếu chỉ chép nguyên văn, lời gọi có thể sai kiểu \texttt{integer}, \texttt{number} hoặc định dạng ngày. Vì vậy, phương pháp trích xuất cần thêm bước chuẩn hóa giá trị (Value Normalizer).

Hệ thống thanh điệu và các biến thể phương ngữ cũng đặt ra yêu cầu đối với bộ truy hồi. Sai dấu hoặc thiếu dấu khi nhập liệu có thể làm thay đổi thực thể địa danh; các cách viết tắt như “TP.HCM”, “HN” hoặc những từ gần nghĩa theo vùng miền như “xe gắn máy”, “xe honda”, “xe máy” cần được xử lý nhất quán trong biểu diễn ngữ nghĩa. Đây là các nguồn biến thiên cần khảo sát khi đánh giá khả năng truy hồi công cụ tiếng Việt.

